% Lösungen Einheit 1 – Nr. 12–23
\einheitenkopf[e1]{Einheit 1 von 4 · Wurzelziehen und Lösbarkeit}

\erg{12}{a) linear \quad b) quadratisch \quad c) quadratisch \quad d) linear \quad e) quadratisch}
\erg{13}{a) zwei \quad b) keine \quad c) zwei \quad d) eine \quad e) keine}
\erg{14}{a) $x_1 = 5$, $x_2 = -5$ \quad b) $x_1 = 9$, $x_2 = -9$ \quad c) $x_1 = 2$, $x_2 = -2$ \quad d) $x_1 = 12$, $x_2 = -12$ \quad e) $x_{1,2} = \pm\sqrt{11} \approx \pm 3{,}32$ \quad f) keine Lösung, ein Quadrat ist nie negativ}
\erg{15}{a) $x^2 = 16$; $x_1 = 4$, $x_2 = -4$ \quad b) $x^2 = 25$; $\pm 5$ \quad c) $x^2 = 9$; $\pm 3$ \quad d) $3x^2 = 27$, $x^2 = 9$; $\pm 3$ \quad e) $4x^2 = 0$, $x^2 = 0$; $x = 0$ (eine Lösung) \quad f) $2x^2 = -8$, $x^2 = -4$; keine Lösung \quad g) $x^2 = 7$; $\pm\sqrt{7} \approx \pm 2{,}65$}
\erg{16}{a) $x - 2 = \pm 3$; $x_1 = 5$, $x_2 = -1$ \quad b) $x_1 = 5$, $x_2 = -3$ \quad c) $x_1 = 7$, $x_2 = 3$ \quad d) $x_1 = 8$, $x_2 = -2$ \quad e) $x + 3 = \pm 7$; $x_1 = 4$, $x_2 = -10$ \quad f) $x + 6 = 0$; $x = -6$ (eine Lösung) \quad g) $x = 2 \pm\sqrt{5}$; $x_1 \approx 4{,}24$, $x_2 \approx -0{,}24$}
\erg{17}{a) (wA) \quad b) $16 + 16 = 32$, (fA) \quad c) $(-6)^2 = 36$, (wA) \quad d) $8^2 = 64$, (fA) \quad e) $2\cdot 4 + 1 = 9$, (wA) \quad f) $(-2)^2 = 4$, (fA)}
\erg{18}{a) 2 \quad b) 1 (Scheitel liegt auf der Geraden) \quad c) 0 (Scheitel liegt über der Geraden) \quad d) 2 (Nullstellen) \quad e) $x_1 = -1$, $x_2 = 3$}
\erg{19}{zum Beispiel: a) $x^2 = -1$ \quad b) $(x - 1)^2 = -4$ \quad c) $2x^2 + 5 = 1$ \quad jede Gleichung, bei der nach dem Freistellen rechts eine negative Zahl steht}
\erg{20}{a) wahr, $x - 5 = 0$, nur $x = 5$ \quad b) falsch, $x + 6 = \pm 4$ ergibt $x_1 = -2$ und $x_2 = -10$ \quad c) jede negative Zahl, zum Beispiel $(x + 6)^2 = -1$}
\erg{21}{a) Die negative Lösung fehlt: $x_1 = 9$, $x_2 = -9$. \quad b) 2 falsch hinübergebracht (addiert statt subtrahiert): $x_1 = 5$, $x_2 = -9$. \quad c) Aus $-9$ gibt es keine Wurzel; $(-3)^2 = 9$, nicht $-9$. Die Gleichung hat keine Lösung.}
\erg{22}{a) $\sqrt{20}$ und $-\sqrt{20}$ ergeben quadriert beide 20. \quad b) Ein Quadrat ist nie negativ, keine Zahl ergibt quadriert $-20$. \quad c) Stimmt nicht: $x - 4 = 0$ hat die Lösung $x = 4$, die Gleichung hat genau eine Lösung.}
\erg{23}{a) $x^2 = 64$, $x = 8$ (die negative Lösung ist keine Länge); die Seite ist 8\,m lang. \quad b) $300\,\ell = 300\,000\,\text{cm}^3$; $r^2 = 300\,000 : (\pi\cdot 90) \approx 1061{,}0$; $r \approx 32{,}6\,\text{cm}$}
